意味をなす句の連なりを文と呼ぶ。
文末の句が主要部となる。

\paragraph{構造}
文は以下の3種類の構造を持つ

\subparagraph{名詞文}
名詞句のみからなる文を名詞文と呼ぶ。
主要部は名詞句となる。

% TODO 例文
\begin{exe}
    \ex \gll @@@\\
    @@@ \\
    \glt @@@
\end{exe}

\subparagraph{コピュラ文}
２つの名詞句がコピュラ"ai"で結ばれた文をコピュラ文と呼ぶ。
主要部は文末の名詞句となる。

% TODO 人名について考察
% 句構造を表す方法
\begin{exe}
    \ex \gll [u\'ak\'a ap\'a piku\'o] ai [XXX] \\
    父の名前 XXX \\
    \glt 父の名前はXXXだ。
\end{exe}

\subparagraph{動詞文}
複数の副詞句、名詞句の後ろに動詞句が並んだ文を動詞句と呼ぶ。
主要部は動詞句となる。

\begin{exe}
    \ex \gll [reap\'a] [t\'iv\'o-p\'a] \\
    私 歩いた \\
    \glt 私は歩いた。
\end{exe}